
\section{The carrier calculus and choice-independent invariants}
\label{sec:carriers}

\subsection{Composition}

\begin{theorem}[Carrier composition]
\label{thm:carrier-composition}
If $C:K\multimap L$ and $D:L\multimap M$ are strong carriers, then $D\circ C:K\multimap M$ is a strong carrier.  Composition is associative, and singleton identity clouds are identities.
\end{theorem}

\begin{proof}
For every source vertex $v$, $C(v)$ is nonempty and every $D(w)$ with $w\in C(v)$ is nonempty.  Therefore their finite union $(D\circ C)(v)$ is nonempty.

Let $\sigma$ be a simplex of $K$.  By the carrier property for $C$,
\[
  \tau=C(\sigma)=\bigcup_{v\in\sigma}C(v)
\]
is a simplex of $L$.  Apply the carrier property for $D$ to $\tau$:
\[
  D(\tau)=\bigcup_{w\in\tau}D(w)
\]
is a simplex of $M$.  Reassociating finite unions gives
\[
  D(\tau)
  =\bigcup_{v\in\sigma}\bigcup_{w\in C(v)}D(w)
  =(D\circ C)(\sigma).
\]
Hence the composite is a carrier.  Associativity and identity follow from associativity and identity of finite union.  Every equality can be proved bit-by-bit by bounded induction.
\end{proof}

\subsection{Selections and contiguity}

\begin{definition}[Selection]
A selection of $C:K\multimap L$ is a vertex map $f:V(K)\to V(L)$ with $f(v)\in C(v)$ for all $v$.
\end{definition}

The least member of each nonempty finite cloud gives a canonical primitive-recursive selection.  This selection is used for execution; correctness does not depend on it.

\begin{theorem}[Selection theorem]
\label{thm:selection}
Every selection $f$ of a strong carrier $C:K\multimap L$ is simplicial.  Any two selections $f,g$ are contiguous: for every simplex $\sigma$ of $K$, the vertex set
\[
  f(\sigma)\cup g(\sigma)
\]
is a simplex of $L$.
\end{theorem}

\begin{proof}
For each $v\in\sigma$, $f(v)\in C(v)$.  Hence
\[
  f(\sigma)\subseteq \bigcup_{v\in\sigma}C(v)=C(\sigma).
\]
The right side is a simplex by the carrier condition, and every subset of a simplex is a simplex.  Thus $f$ is simplicial.

The same argument applied simultaneously to $f$ and $g$ gives
\[
  f(\sigma)\cup g(\sigma)\subseteq C(\sigma),
\]
so the union is a simplex.  This is exactly contiguity.
\end{proof}

\subsection{Chains over $\Ftwo$}

For an ordered finite complex $K$, let $C_q(K;\Ftwo)$ be the finite-dimensional vector space with basis the $q$-simplices.  The boundary is
\[
  \partial[v_0,\ldots,v_q]
  =\sum_{i=0}^{q}[v_0,\ldots,\widehat{v_i},\ldots,v_q]
\]
with coefficients modulo $2$.  Each $(q-2)$-face appears exactly twice in $\partial^2$, so $\partial^2=0$.

A simplicial vertex map $f$ sends a simplex to the sorted set of image vertices; if two source vertices have the same image, the normalized term is zero.  Extending linearly gives a chain map $f_\#$.

All operations are bounded parity folds over finite simplex tables.  No quotient type is needed inside arithmetic: a chain is simply a bitset indexed by the enumerated simplices.

\subsection{The prism operator}

Let $f,g:K\to L$ be contiguous.  For a basis simplex $\sigma=[v_0,\ldots,v_q]$, define over $\Ftwo$
\[
  P_q(\sigma)=
  \sum_{i=0}^{q}
  [f(v_0),\ldots,f(v_i),g(v_i),\ldots,g(v_q)],
\]
where a term with a repeated vertex is normalized to zero.  Contiguity ensures that every nondegenerate term is a simplex of $L$.

\begin{theorem}[Finite prism identity]
\label{thm:prism}
For contiguous simplicial maps $f,g:K\to L$,
\[
  \partial P+P\partial=f_\#+g_\#
\]
as maps of mod-$2$ chain complexes.
\end{theorem}

\begin{proof}
Fix $\sigma=[v_0,\ldots,v_q]$ and expand $\partial P_q(\sigma)$.  A term of $P_q(\sigma)$ has an $f$-block ending at $v_i$ and a $g$-block beginning at the same $v_i$.  Deleting a vertex strictly inside the $f$-block or strictly inside the $g$-block produces exactly the corresponding term of $P_{q-1}(\partial\sigma)$.  These terms therefore occur once in $\partial P$ and once in $P\partial$ and cancel modulo $2$.

At the seam, deleting $g(v_i)$ from the $i$-th prism term produces the same normalized simplex as deleting $f(v_{i+1})$ from the $(i+1)$-st prism term.  These adjacent seam terms also cancel in pairs.  The two unpaired end terms are obtained by deleting $f(v_0)$ from the first prism and deleting $g(v_q)$ from the last prism.  They are precisely $g_\#(\sigma)$ and $f_\#(\sigma)$.  Thus
\[
  (\partial P+P\partial)(\sigma)=f_\#(\sigma)+g_\#(\sigma).
\]

If a tuple has repeated vertices, normalization maps it to zero.  The cancellation pairing respects this normalization: paired words have the same multiset of vertices, so either both normalize to the same simplex or both vanish.  Therefore the identity survives degeneracies.  Linearity proves it for every chain.

For formalization, the preceding prose becomes a bounded index pairing.  Each emitted term is tagged by $(i,j)$, where $i$ is the prism seam and $j$ is the deleted position.  A primitive-recursive involution pairs every tag except the two endpoints; a parity fold proves cancellation by quantifier-free induction.
\end{proof}

\begin{corollary}[Choice-independent homology map]
\label{cor:choice-independent}
Any two selections of a strong carrier induce the same map
\[
  H_q(K;\Ftwo)\longrightarrow H_q(L;\Ftwo)
\]
for every $q$.
\end{corollary}

\begin{proof}
Let $z$ be a cycle.  By \cref{thm:prism},
\[
  f_\#z+g_\#z=\partial Pz+P\partial z=\partial Pz.
\]
Thus $f_\#z$ and $g_\#z$ differ by a boundary and represent the same homology class.
\end{proof}

\subsection{Functorial tower homology}

Every refinement carrier $R_k:K_{k+1}\multimap K_k$ now induces a canonical homology map $r_k$, because the chosen selection is irrelevant by \cref{cor:choice-independent}.  Composition of carriers induces composition of homology maps: selecting from a composite and composing selections are both selections of a common composite carrier, hence contiguous.

\begin{theorem}[Tower homology]
\label{thm:tower-homology}
A certified observation tower determines, for every $q$, an inverse system of finite-dimensional $\Ftwo$-vector spaces
\[
  \cdots\longrightarrow H_q(K_{k+1};\Ftwo)
  \xrightarrow{r_k}
  H_q(K_k;\Ftwo)
  \longrightarrow\cdots.
\]
A coherent tower map determines a morphism of inverse systems, independent of all finite selections.  Composition of tower maps is functorial.
\end{theorem}

\begin{proof}
Choose least selections for execution.  The refinement maps are well defined on homology by \cref{cor:choice-independent}.  The coherence certificate for a tower map says that the two finite vertex maps around each naturality square are carried by a common carrier.  They are therefore contiguous and induce equal homology maps.  This is exactly commutativity of the square on homology.  Carrier composition represents map composition, and \cref{thm:carrier-composition,cor:choice-independent} prove associativity and independence of all selection choices.
\end{proof}

\subsection{Integer coefficients}

Integer chains require signs.  One can use the signed prism
\[
  P_q(\sigma)=\sum_{i=0}^{q}(-1)^i
  [f(v_0),\ldots,f(v_i),g(v_i),\ldots,g(v_q)]
\]
and the identity $\partial P+P\partial=g_\#-f_\#$.  Signed integers are coded by pairs of naturals or a parity-magnitude code.  Homology groups can be computed by Smith normal form.

For \ipra, the easiest robust interface is a \emph{checked Smith certificate}: matrices $U,D,V$ with $UAV=D$, determinants $\pm1$, diagonal $D$, and the divisibility chain.  A bounded verifier checks the certificate.  A separate bounded Euclidean algorithm may produce it when desired.  The mod-$2$ theory is the minimal core because Gaussian elimination and parity cancellation are especially transparent.

\subsection{Acyclic carriers as an extension}

Strong carriers may be generalized.  Assign to each source simplex $\sigma$ a target subcomplex $A(\sigma)$ with face monotonicity and a supplied finite contraction or acyclicity certificate.  The usual acyclic-carrier induction constructs a chain map and a chain homotopy \cite{Munkres}.  In \ipra, ``acyclic'' must not be an unexamined semantic adjective: the certificate must provide bounded filling operators in the required dimensions.  The strong carrier is the special case where $A(\sigma)$ is one simplex and filling is canonical.
